% !TeX program = pdflatex
\documentclass[aspectratio=169,xcolor={dvipsnames,table}]{beamer}

\input{../../_en_preambulo_beamer}
% Comandos y macros estadísticas compartidas para las presentaciones Beamer en inglés (Metropolis)

% Conjuntos numéricos básicos
\newcommand{\R}{\mathbb{R}}
\newcommand{\N}{\mathbb{N}}
\newcommand{\Q}{\mathbb{Q}}
\newcommand{\Z}{\mathbb{Z}}
\newcommand{\set}[1]{\left\{ #1 \right\}}
\newcommand{\abs}[1]{\left|#1\right|}
\newcommand{\norm}[1]{\left\| #1 \right\|}

% Comandos estadísticos y probabilísticos del curso
\newcommand{\Var}{\operatorname{Var}}
\newcommand{\cov}{\operatorname{Cov}}
\newcommand{\corr}{\rho}
\newcommand{\comb}[2]{\begin{pmatrix} #1 \\ #2 \end{pmatrix}}
\newcommand{\s}{\sigma}
\newcommand{\particion}{\mathfrak{P}}
\newcommand{\card}[1]{\#\left( #1 \right)}
\newcommand{\seq}[3]{\set{#1}_{#2}^{#3}}
\newcommand{\E}{\mathbb{E}}
\newcommand{\Prob}{\mathbb{P}}

% Entornos pedagógicos y cajas en inglés académico natural (soportando tanto nombres en inglés como en español)
\newenvironment{discussion}{%
  \begin{alertblock}{Discussion Question}%
}{%
  \end{alertblock}%
}
\newenvironment{discusion}{%
  \begin{alertblock}{Discussion Question}%
}{%
  \end{alertblock}%
}

\newenvironment{reflection}{%
  \begin{alertblock}{Classroom Reflection}%
}{%
  \end{alertblock}%
}
\newenvironment{reflexion}{%
  \begin{alertblock}{Classroom Reflection}%
}{%
  \end{alertblock}%
}

\newenvironment{paradox}{%
  \begin{alertblock}{Exploratory Paradox}%
}{%
  \end{alertblock}%
}
\newenvironment{paradoja}{%
  \begin{alertblock}{Exploratory Paradox}%
}{%
  \end{alertblock}%
}

\newenvironment{motivation}{%
  \begin{exampleblock}{Motivation: Data Science in Practice}%
}{%
  \end{exampleblock}%
}
\newenvironment{motivacion}{%
  \begin{exampleblock}{Motivation: Data Science in Practice}%
}{%
  \end{exampleblock}%
}

\newenvironment{quickchallenge}{%
  \begin{block}{Quick Team Challenge}%
}{%
  \end{block}%
}
\newenvironment{retoclase}{%
  \begin{block}{Quick Team Challenge}%
}{%
  \end{block}%
}


\title{Continuous Expectation and Variance}
\subtitle{Section 04.03 --- Integral Definition, LOTUS, and Variance Decomposition}
\author[J. Castillo Colmenares]{Juliho Castillo Colmenares}
\institute[Tec de Monterrey]{Tecnologico de Monterrey}
\date{\vspace{-1.5cm}}

\begin{document}

% SLIDE 01: Title Page
\begin{frame}[plain]
	\titlepage
\end{frame}

% SLIDE 02: Roadmap
\begin{frame}{Roadmap --- Chapter 04: Continuous Random Variables}
	\vspace{-0.15cm}
	\begin{block}{Curricular Structure of Unit 3}
		\footnotesize
		\begin{enumerate}\itemsep=0.03cm
			\item \textcolor{gray}{Section 04.01: PDF and Continuous Support}
			\item \textcolor{gray}{Section 04.02: Continuous CDF and Quantiles}
			\item \textbf{\textcolor{TecRojo}{Section 04.03: Expectation, Variance, and Continuous LOTUS}} $\leftarrow$ \emph{Current focus}
			\item \textcolor{gray}{Section 04.04: Continuous Uniform Distribution}
			\item \textcolor{gray}{Section 04.05: Normal Distribution and Z-Score}
			\item \textcolor{gray}{Section 04.06: Exponential Distribution and Memoryless Processes}
			\item \textcolor{gray}{Section 04.07: Gamma, Beta, and Weibull Distributions}
		\end{enumerate}
	\end{block}
	\vspace{-0.2cm}
	\begin{alertblock}{Learning Objective}
		\scriptsize
		Define continuous expectation $\E[X] = \int x f(x)\,dx$ and variance $\Var(X) = \int (x-\mu)^{2} f(x)\,dx$, prove the LOTUS theorem for transformations, and apply the law of total variance to hierarchical models.
	\end{alertblock}
\end{frame}

% SLIDE 03: Motivation
\begin{frame}{From PDF to Moments: Integration as a Natural Operation}
	\footnotesize
	\begin{motivacion}[Why continuous expectation?]
		We know the PDF and CDF of a continuous variable. Now we need numerical summaries that capture the essential properties: the \emph{expectation} $\E[X]$ (center of mass of the distribution) and the \emph{variance} $\Var(X)$ (dispersion around the center). These moments guide decisions in inference, simulation, and risk analysis.
	\end{motivacion}
	\vspace{0.15cm}
	\pause
	\begin{itemize}\itemsep=0.1cm
		\item \textbf{Expectation:} $\E[X] = \int x f(x)\,dx$, a natural generalization of $\E[X] = \sum x P(X=x)$ via Riemann sums. \pause
		\item \textbf{Variance:} $\Var(X) = \E[(X - \mu)^{2}] = \E[X^{2}] - (\E[X])^{2}$, measures mean squared dispersion. \pause
		\item \textbf{LOTUS:} $\E[g(X)] = \int g(x) f(x)\,dx$ avoids deriving the distribution of $g(X)$.
	\end{itemize}
\end{frame}

% SLIDE 03B: Definition and Properties
\begin{frame}{Definition of Continuous Expectation and Variance}
	\vspace{-0.15cm}
	\begin{columns}[T]
		\begin{column}{0.48\textwidth}
			\begin{block}{Expectation}
				\footnotesize
				$\E[X] = \int_{-\infty}^{\infty} x f(x)\,dx$ under convergence $\int |x| f(x)\,dx < \infty$. For light-tailed distributions, the integral converges; for Cauchy, it diverges.
			\end{block}
		\end{column}
		\begin{column}{0.48\textwidth}
			\begin{alertblock}{Variance and LOTUS}
				\footnotesize
				$\Var(X) = \E[(X-\mu)^{2}] = \E[X^{2}] - (\E[X])^{2}$. LOTUS: $\E[g(X)] = \int g(x) f(x)\,dx$ without deriving the distribution of $g(X)$. Linearity: $\E[aX+b] = a\E[X]+b$.
			\end{alertblock}
		\end{column}
	\end{columns}
\end{frame}

% SLIDE 03C: Continuous Variance
\begin{frame}{Continuous Variance and Second Moments}
\scriptsize
\begin{block}{Definition}
The variance of a continuous variable is
\[
\Var(X)=\E[(X-\mu)^2]=\int (x-\mu)^2f_X(x)\,dx
       =\E[X^2]-\mu^2.
\]
\end{block}
\begin{alertblock}{Interpretation}
Variance measures quadratic dispersion around $\mu=\E[X]$ and has squared units;
the standard deviation $\sigma=\sqrt{\Var(X)}$ returns to the original scale.
\end{alertblock}
\end{frame}

% SLIDE 03D: Linearity and Variance Properties
\begin{frame}{Linearity of Expectation and Variance Properties}
	\vspace{-0.15cm}
	\begin{columns}[T]
		\begin{column}{0.48\textwidth}
			\begin{block}{Variance Properties}
				\footnotesize
				\begin{itemize}\itemsep=0.03cm
					\item $\Var(X) \ge 0$ always, with equality iff $X$ is constant.
					\item $\Var(aX + b) = a^{2} \Var(X)$ (affine linearity).
				\end{itemize}
			\end{block}
		\end{column}
		\begin{column}{0.48\textwidth}
			\begin{alertblock}{Linearity Theorem}
				\footnotesize
				$\E[aX+b] = a\E[X]+b$ and $\E[X+Y] = \E[X]+\E[Y]$, even for \emph{dependent} $X, Y$. In contrast, $\Var(X+Y) = \Var(X) + \Var(Y) + 2\cov(X,Y)$ requires independence to drop the covariance term.
			\end{alertblock}
		\end{column}
	\end{columns}
\end{frame}

% SLIDE 03E: Linearity of Expectation
\begin{frame}{Linearity of Continuous Expectation}
\scriptsize
For constants $a,b$ and integrable variables,
\[
\E[aX+b]=a\E[X]+b,\qquad \Var(aX+b)=a^2\Var(X).
\]
Linearity holds without independence; independence is needed only for simplifying
variances of sums through covariance terms.
\end{frame}

% SLIDE 04: Central Moments and Law of Total Variance
\begin{frame}{Central Moments and Law of Total Variance}
	\vspace{-0.15cm}
	\begin{columns}[T]
		\begin{column}{0.48\textwidth}
			\begin{block}{Skewness and Kurtosis}
				\footnotesize
				Skewness: $\gamma_{1} = \E[(X-\mu)^{3}]/\sigma^{3}$ measures asymmetry. Kurtosis: $\gamma_{2} = \E[(X-\mu)^{4}]/\sigma^{4} - 3$ measures tail heaviness. For Exponential(1): $\gamma_{1} = 2$, $\gamma_{2} = 6$.
			\end{block}
		\end{column}
		\begin{column}{0.48\textwidth}
			\begin{alertblock}{Law of Total Variance}
				\footnotesize
				$\Var(X) = \E[\Var(X \mid Y)] + \Var(\E[X \mid Y])$. Fundamental in ANOVA: decomposes total variance into within-group and between-group components. For hierarchical models with $Y$ as group index.
			\end{alertblock}
		\end{column}
	\end{columns}
\end{frame}

% SLIDE 05: Applications
\begin{frame}{Applications in Engineering, Science, and Data Science}
	\vspace{-0.1cm}
	\begin{itemize}\itemsep=0.1cm
		\item \textbf{Uncertainty Propagation:} If $X = \mu + \epsilon$ with $\epsilon \sim N(0, \sigma^2)$, the standard error of the mean of $n$ measurements is $\sigma/\sqrt{n}$. \pause
		\item \textbf{Risk Quantification:} VaR at 95\% is the quantile $q_{0.05}$ of the loss distribution; CVaR is $\E[X \mid X \le q_{0.05}]$. \pause
		\item \textbf{Regression Analysis:} Residuals must have zero mean and constant variance (homoscedasticity). \pause
		\item \textbf{Quality Control:} $\bar{x}$ charts use $\mu \pm 3\sigma/\sqrt{n}$ as control limits.
	\end{itemize}
\end{frame}

% SLIDE 06: Python Lab Block 1
\begin{frame}[fragile]{Python Lab: Expectation and LOTUS}
	\vspace{-0.05cm}
	\scriptsize
	Computation of $\E[X]$, $\E[X^2]$ and LOTUS verification for $g(X) = \sqrt{X}$ and $g(X) = \log(X)$ (\texttt{04.03\_expectation\_and\_variance.py}):
	\vspace{0.02cm}
	{\lstset{basicstyle=\fontsize{4.5pt}{5.5pt}\selectfont\ttfamily}
	\lstinputlisting[firstline=15, lastline=38]{../../code/04_variables_aleatorias_continuas/04.03_expectation_and_variance.py}}
\end{frame}

% SLIDE 07: Python Lab Block 2
\begin{frame}[fragile]{Python Lab: Variance and Central Moments}
	\vspace{-0.1cm}
	\scriptsize
	Variance and central moment computation (skewness, kurtosis) for Uniform, Exponential, Normal, and law of total variance demonstration:
	\vspace{0.0cm}
	{\lstset{basicstyle=\fontsize{4pt}{5pt}\selectfont\ttfamily}
	\lstinputlisting[firstline=45, lastline=72]{../../code/04_variables_aleatorias_continuas/04.03_expectation_and_variance.py}}
\end{frame}

% SLIDE 08: Python Lab Block 3
\begin{frame}[fragile]{Python Lab: Monte Carlo and Error Propagation}
	\vspace{-0.05cm}
	\scriptsize
	Monte Carlo simulation with $N = 250{,}000$ and error propagation analysis in measurement:
	\vspace{0.02cm}
	{\lstset{basicstyle=\fontsize{4.5pt}{5.5pt}\selectfont\ttfamily}
	\lstinputlisting[firstline=85, lastline=108]{../../code/04_variables_aleatorias_continuas/04.03_expectation_and_variance.py}}
\end{frame}

% SLIDE 09: Python Lab Terminal Output
\begin{frame}[fragile]{Python Lab: Terminal Output}
	\vspace{-0.1cm}
	\begin{block}{}
		\fontsize{4pt}{5pt}\selectfont
		\begin{verbatim}
=== Block 1: Expectation & LOTUS ===
Triangular(0,1): E[X]=0.666667 | E[X^2]=0.500000 | Var=0.055556
LOTUS E[sqrt(X)]=0.800000 | LOTUS E[log(X)]=-0.499993
Linearity E[3.5X+10]=12.333333

=== Block 2: Variance & Central Moments ===
Uniform(0,1): E[X]=0.500000 | Var=0.083333
Exponential(1.0): Skewness=2.000000 | Excess kurtosis=6.000000
Hierarchical: total Var=4.2187 = within 1.0221 + between 3.2991

=== Block 3: Monte Carlo & Error Propagation ===
Exponential MC (N=250,000): mean=1.0011 | var=1.0003
  Skewness=1.9909 | Kurtosis=5.8731
Error propagation (sigma=0.1, true=5.0):
  n=1024: SE=0.0031 | CI half-width=0.0061
  Required n for SE<0.01: 100
		\end{verbatim}
	\end{block}
\end{frame}









% SLIDE 18: Closing and Next Steps
\begin{frame}{Closing Section and Modular Perspective}
	\vspace{-0.2cm}
	\begin{block}{Synthesis: Moments as Distributional Summary}
		\scriptsize
		Continuous expectation $\E[X] = \int x f(x)\,dx$ and variance $\Var(X) = \int (x - \mu)^{2} f(x)\,dx$ summarize the essential properties of any distribution. The LOTUS theorem extends these ideas to transformations without deriving new PDFs, and the law of total variance connects moments with hierarchical structures.
	\end{block}
	\vspace{0.05cm}
	\pause
	\begin{alertblock}{Key Lessons and Next Step}
		\begin{itemize}\itemsep=0.05cm
			\item \textbf{Expectation:} $\E[X]$ is the center of mass; the integral $\int x f(x)\,dx$ generalizes discrete summation.
			\item \textbf{Variance:} $\Var(X) = \E[X^{2}] - (\E[X])^{2}$; variances sum only for independent variables.
			\item \textbf{LOTUS:} $\E[g(X)] = \int g(x) f(x)\,dx$ avoids deriving the distribution of $g(X)$.
			\item \textbf{Next Section:} \textcolor{TecRojo}{Continuous Uniform Distribution}.
		\end{itemize}
	\end{alertblock}
\end{frame}

% SLIDE 20: Modular Perspective
\begin{frame}{Modular Perspective --- The Next Challenge}
\scriptsize
\begin{alertblock}{Next Session: Continuous Distributions}
Moments summarize a distribution; the next sections apply them to Uniform,
Normal, Gamma, and other continuous families used in engineering and data science.
\end{alertblock}
\begin{block}{Recommended Preparation}
Review conditional expectation and the law of total variance before applying these
identities to sampling distributions.
\end{block}
\end{frame}

\end{document}
